\documentclass[12pt,a4paper]{article}

% ============================================================
% PACKAGES
% ============================================================
\usepackage[utf8]{inputenc}
\usepackage[T1]{fontenc}
\usepackage{lmodern}
\usepackage[french]{babel}
\usepackage{geometry}
\usepackage{graphicx}
\usepackage[table]{xcolor}
\usepackage{booktabs}
\usepackage{array}
\usepackage{fancyhdr}
\usepackage{titlesec}
\usepackage{enumitem}
\usepackage{mdframed}
\usepackage{float}
\usepackage{microtype}
\usepackage{setspace}
\usepackage{hyperref}

% ============================================================
% MISE EN PAGE
% ============================================================
\geometry{
  headheight=14pt,
  top=2.6cm,
  bottom=2.8cm,
  left=2.8cm,
  right=2.8cm
}
\setstretch{1.12}

% ============================================================
% COULEURS
% ============================================================
\definecolor{primary}{RGB}{28, 90, 168}
\definecolor{primarylight}{RGB}{88, 145, 220}
\definecolor{secondary}{RGB}{232, 240, 252}
\definecolor{success}{RGB}{34, 130, 90}
\definecolor{warning}{RGB}{200, 120, 40}
\definecolor{darkgray}{RGB}{45, 55, 72}
\definecolor{lightgray}{RGB}{95, 110, 140}

% ============================================================
% EN-TETE ET PIED DE PAGE
% ============================================================
\pagestyle{fancy}
\fancyhf{}
\fancyhead[L]{\small\sffamily\textcolor{primary!85!black}{Projet M1 GIL --- Gestion de flotte}}
\fancyhead[R]{\small\sffamily\textcolor{primary!85!black}{Univ. Rouen Normandie \textperiodcentered{} 2025--2026}}
\fancyfoot[C]{\small\sffamily\textcolor{primary!70!black}{\thepage}}
\renewcommand{\headrulewidth}{0.4pt}
\renewcommand{\headrule}{\hbox to\headwidth{\color{primary!35}\leaders\hrule height \headrulewidth\hfill}}
\renewcommand{\footrulewidth}{0pt}

% ============================================================
% STYLE DES TITRES
% ============================================================
\titleformat{\section}
  {\Large\sffamily\bfseries\color{primary}}
  {\thesection.}{0.55em}{}

\titleformat{\subsection}
  {\large\sffamily\bfseries\color{primary!85!black}}
  {\thesubsection.}{0.45em}{}

\titleformat{\subsubsection}
  {\normalsize\sffamily\bfseries\color{primary!75!black}}
  {\thesubsubsection.}{0.4em}{}

\titlespacing*{\section}{0pt}{1.6ex plus .2ex}{1ex plus .1ex}
\titlespacing*{\subsection}{0pt}{1.3ex plus .15ex}{0.7ex plus .1ex}

% ============================================================
% BOITE INFO
% ============================================================
\newmdenv[
  backgroundcolor=secondary,
  linecolor=primary,
  linewidth=1.1pt,
  leftline=true,
  rightline=false,
  topline=false,
  bottomline=false,
  innerleftmargin=11pt,
  innerrightmargin=11pt,
  innertopmargin=9pt,
  innerbottommargin=9pt,
  skipabove=0.6em,
  skipbelow=0.6em
]{infobox}

\hypersetup{
  colorlinks=true,
  linkcolor=primary!90!black,
  urlcolor=primarylight,
  pdfborderstyle={/S/U/W 0}
}

% ============================================================
% DOCUMENT
% ============================================================
\begin{document}

% ------------------------------------------------------------
% PAGE DE TITRE
% ------------------------------------------------------------
\begin{titlepage}
  \centering
  \vspace*{1cm}

  {\Large\textbf{Universit\'{e} de Rouen Normandie}\\[0.3cm]
  \large Master 1 -- G\'{e}nie Informatique et Logiciel}

  \vspace{1.5cm}
  {\color{primary!55}\rule{\linewidth}{0.5mm}}\\[0.45cm]

  {\Huge\sffamily\bfseries\color{primary} Projet Gestion de Flotte\\[0.35cm]}
  {\LARGE\sffamily\bfseries\color{primary!80!black} Rapport de Semaine 5}\\[0.35cm]
  {\large\sffamily\color{primary!65!black} Service \'{E}v\'{e}nements \& Alertes --- Service Localisation}

  {\color{primary!55}\rule{\linewidth}{0.5mm}}\\[1.5cm]

  \begin{minipage}{0.45\textwidth}
    \begin{flushleft}
      \textbf{\'{E}tudiants :}\\
      Ahcene Amouchas \\
      Florian Pepin \\[0.5cm]
      \textbf{Encadrants :}\\
      Lydia \& Luc
    \end{flushleft}
  \end{minipage}
  \hfill
  \begin{minipage}{0.45\textwidth}
    \begin{flushright}
      \textbf{Ann\'{e}e universitaire :}\\
      2025 -- 2026 \\[0.5cm]
      \textbf{Semaine :}\\
      5 / 7
    \end{flushright}
  \end{minipage}

  \vfill
  {\large 9 avril 2026}
\end{titlepage}

\tableofcontents
\newpage

% ============================================================
\section{Introduction}
% ============================================================

La semaine~5 ajoute deux nouveaux microservices au syst\`{e}me : le \textbf{service \'{e}v\'{e}nements}, qui centralise toutes les alertes de la flotte, et le \textbf{service localisation}, qui g\`{e}re les positions GPS des v\'{e}hicules en temps r\'{e}el.

\subsection{R\'{e}partition du travail}

\begin{itemize}[leftmargin=*]
  \item \textbf{Florian Pepin} --- Service \'{e}v\'{e}nements \& alertes : consommation Kafka, API REST, int\'{e}gration GraphQL, tests, d\'{e}ploiement.
  \item \textbf{Ahcene Amouchas} --- Service localisation : ingestion GPS via gRPC, stockage s\'{e}ries temporelles, g\'{e}ofencing, simulation.
\end{itemize}

Les deux services ont \'{e}t\'{e} d\'{e}velopp\'{e}s en parall\`{e}le avec un contrat d'interface Kafka pr\'{e}d\'{e}fini, permettant \`{a} chacun d'avancer ind\'{e}pendamment.

% ============================================================
\newpage
\section{Service \'{E}v\'{e}nements \& Alertes}
% ============================================================

\begin{infobox}
\textbf{R\^{o}le :} Agr\'{e}gateur central passif. Le service \'{e}coute les topics Kafka des autres microservices, transforme chaque \'{e}v\'{e}nement pertinent en alerte persistante, et expose une API pour que les op\'{e}rateurs puissent les consulter et les g\'{e}rer.
\end{infobox}

\subsection{Mod\`{e}le d'alerte}

Chaque alerte stocke un \textbf{type}, une \textbf{s\'{e}v\'{e}rit\'{e}} et un \textbf{statut}, ainsi que des r\'{e}f\'{e}rences logiques (UUID) vers le v\'{e}hicule et le conducteur concern\'{e}s --- sans cl\'{e} \'{e}trang\`{e}re inter-services.

\begin{table}[H]
\centering
\small
\begin{tabular}{|l|l|l|}
  \hline
  \rowcolor{primary!15}
  \textbf{Type d'alerte} & \textbf{Source} & \textbf{S\'{e}v\'{e}rit\'{e}} \\
  \hline
  \texttt{MAINTENANCE\_OVERDUE}    & Service Maintenance (Kafka) & WARNING / INFO \\
  \texttt{LICENSE\_EXPIRING}       & Service Conducteurs (Kafka) & WARNING \\
  \texttt{LICENSE\_EXPIRED}        & Service Conducteurs (Kafka) & CRITICAL \\
  \texttt{GEOFENCING\_BREACH}      & Service Localisation (Kafka) & HIGH \\
  \texttt{SPEED\_EXCEEDED}         & Service Localisation (Kafka) & WARNING \`{a} CRITICAL \\
  \texttt{VEHICLE\_IMMOBILIZED}    & Service Localisation (Kafka) & WARNING \\
  \hline
\end{tabular}
\caption{Types d'alertes support\'{e}s et leur origine}
\end{table}

Le cycle de vie d'une alerte suit trois \'{e}tats : \textbf{ACTIVE} \textrightarrow{} \textbf{ACKNOWLEDGED} \textrightarrow{} \textbf{RESOLVED}. Une alerte peut aussi \^{e}tre r\'{e}solue directement sans passer par l'\'{e}tat acquitt\'{e}.

\subsection{Consommation Kafka}

Le service s'abonne \`{a} trois topics :

\begin{itemize}[leftmargin=*]
  \item \textbf{\texttt{flotte.alertes.events}} : messages publi\'{e}s par le service maintenance lors d'un retard de maintenance. Chaque message produit une alerte \texttt{MAINTENANCE\_OVERDUE}.

  \item \textbf{\texttt{flotte.conducteurs.events}} : messages publi\'{e}s par le service conducteurs. Seuls les \'{e}v\'{e}nements \texttt{LICENSE\_EXPIRING} et \texttt{LICENSE\_EXPIRED} g\'{e}n\`{e}rent une alerte ; les autres (\texttt{DRIVER\_CREATED}, etc.) sont ignor\'{e}s silencieusement.

  \item \textbf{\texttt{flotte.location.events}} : consumer d\'{e}j\`{a} d\'{e}ploy\'{e} et pr\^{e}t \`{a} recevoir les \'{e}v\'{e}nements du service localisation (\texttt{GEOFENCING\_BREACH}, \texttt{SPEED\_EXCEEDED}\ldots). Il attend silencieusement les premiers messages sans erreur.
\end{itemize}

Un choix technique notable : tous les messages sont d\'{e}s\'{e}rialis\'{e}s manuellement (d\'{e}serialiseur String brut + Jackson), ce qui rend le service compatible avec n'importe quel producteur quel que soit son langage (Java, Node.js, etc.) sans d\'{e}pendance aux en-t\^{e}tes propri\'{e}taires de Spring Kafka.

\subsection{API REST}

L'ensemble des routes est expos\'{e} sous le pr\'{e}fixe \texttt{/api/alerts} :

\begin{table}[H]
\centering
\small
\begin{tabular}{|l|l|l|}
  \hline
  \rowcolor{primary!15}
  \textbf{M\'{e}thode} & \textbf{Endpoint} & \textbf{Description} \\
  \hline
  GET   & \texttt{/api/alerts}                     & Liste avec filtres (statut, s\'{e}v\'{e}rit\'{e}, type, vehicleId\ldots) \\
  GET   & \texttt{/api/alerts/\{id\}}              & D\'{e}tail d'une alerte \\
  GET   & \texttt{/api/alerts/stats}               & Nombre d'alertes actives \\
  PATCH & \texttt{/api/alerts/\{id\}/acknowledge}  & Acquitter une alerte (admin, manager) \\
  PATCH & \texttt{/api/alerts/\{id\}/resolve}      & R\'{e}soudre une alerte (admin, manager, technicien) \\
  \hline
\end{tabular}
\caption{Endpoints REST du service \'{e}v\'{e}nements}
\end{table}

L'identifiant de l'op\'{e}rateur qui acquitte une alerte est extrait directement du jeton JWT Keycloak (claim \texttt{sub}), sans le passer dans le corps de la requ\^{e}te. Cela \'{e}vite tout risque d'usurpation d'identit\'{e}.

\subsection{Int\'{e}gration GraphQL}

Les op\'{e}rations \texttt{alerts}, \texttt{alert}, \texttt{acknowledgeAlert} et \texttt{resolveAlert}, jusqu'ici en \'{e}tat \texttt{NOT\_IMPLEMENTED} dans la gateway, sont d\'{e}sormais branch\'{e}es sur ce service.

Un point notable : lorsqu'une alerte est r\'{e}cup\'{e}r\'{e}e via GraphQL, les champs imbriqués \texttt{vehicle} et \texttt{driver} sont r\'{e}solus automatiquement par la gateway en interrogeant les services correspondants en parall\`{e}le. Si l'un des services est indisponible, le champ renvoie \texttt{null} sans faire \'{e}chouer la requ\^{e}te.

\begin{table}[H]
\centering
\small
\begin{tabular}{|l|c|}
  \hline
  \rowcolor{primary!15}
  \textbf{Op\'{e}ration GraphQL} & \textbf{Statut} \\
  \hline
  \texttt{alerts(\ldots)}            & \textcolor{success}{\textbf{OK}} \\
  \texttt{alert(id)}                 & \textcolor{success}{\textbf{OK}} \\
  \texttt{acknowledgeAlert(id)}      & \textcolor{success}{\textbf{OK}} \\
  \texttt{resolveAlert(id)}          & \textcolor{success}{\textbf{OK}} \\
  \texttt{Alert.vehicle} (r\'{e}solution crois\'{e}e) & \textcolor{success}{\textbf{OK}} \\
  \texttt{Alert.driver} (r\'{e}solution crois\'{e}e)  & \textcolor{success}{\textbf{OK}} \\
  Subscriptions WebSocket            & \textcolor{warning}{\textbf{Semaine 7}} \\
  \hline
\end{tabular}
\caption{Op\'{e}rations GraphQL du domaine alertes}
\end{table}

\subsection{Tests}

\begin{itemize}[leftmargin=*]
  \item \textbf{9 tests unitaires} (JUnit~5 + Mockito) couvrant la logique m\'{e}tier de l'\texttt{AlertService} : cr\'{e}ation, acquittement, r\'{e}solution, idempotence, cas d'erreur.
  \item \textbf{0 \'{e}chec}, couverture $>$ 80\% sur la couche service.
  \item \textbf{Collection Bruno} : 9 requ\^{e}tes REST et 6 requ\^{e}tes GraphQL, toutes valides en environnement Minikube.
\end{itemize}

\subsection{D\'{e}ploiement}

Le service est disponible dans les deux environnements cibles :

\begin{itemize}[leftmargin=*]
  \item \textbf{Docker Compose} : conteneur \texttt{flotte-events} sur le port 8083, d\'{e}pendant de PostgreSQL (healthcheck), Kafka et Keycloak.
  \item \textbf{Kubernetes / Helm} : ajout dans la map \texttt{services} du chart \texttt{fleet-app} (Deployment, Service, probes de sant\'{e}), cr\'{e}ation de \texttt{events\_db} dans le script \texttt{initdb} du chart \texttt{fleet-infra}, route \texttt{/api/alerts} ajout\'{e}e \`{a} l'Ingress.
\end{itemize}

% ============================================================
\newpage
\section{Service Localisation}
% ============================================================

\begin{infobox}
\textbf{R\^{o}le :} Recevoir et stocker les positions GPS des v\'{e}hicules en continu, d\'{e}tecter les franchissements de zones (g\'{e}ofencing), et publier les alertes correspondantes sur Kafka.
\end{infobox}

\subsection{Choix technologiques}

Contrairement aux autres microservices (Java / Spring Boot), le service localisation est d\'{e}velopp\'{e} en \textbf{Node.js avec NestJS (TypeScript)}. Ce choix se justifie par la nature du probl\`{e}me : l'ingestion GPS implique un tr\`{e}s grand nombre de connexions simultan\'{e}es \`{a} faible charge unitaire, ce que le mod\`{e}le \'{e}v\'{e}nementiel non-bloquant de Node.js g\`{e}re bien mieux qu'un mod\`{e}le thread-par-requ\^{e}te.

La communication entre les boitiers embarqu\'{e}s et le service utilise \textbf{gRPC} (Protocol Buffers) plut\^{o}t que HTTP/REST. gRPC permet d'\'{e}tablir un tunnel persistant entre le v\'{e}hicule et le serveur, \'{e}vitant de rouvrir une connexion TCP \`{a} chaque point GPS envoy\'{e} --- ce qui est co\^{u}teux \`{a} haute fr\'{e}quence.

Pour la persistance, \textbf{PostgreSQL} est utilis\'{e} avec l'extension \textbf{TimescaleDB}, qui convertit la table de positions en \textit{hypertable} (partitionn\'{e}e automatiquement par plage temporelle). Cela garantit des performances stables m\^{e}me avec des millions de points stock\'{e}s.

\subsection{Contrat gRPC}

Le fichier \texttt{location.proto} d\'{e}finit deux m\'{e}thodes principales :

\begin{itemize}[leftmargin=*]
  \item \textbf{\texttt{StreamPositions}} : streaming bidirectionnel. Le v\'{e}hicule pousse ses coordonn\'{e}es GPS (latitude, longitude, vitesse, cap) en continu dans un tunnel persistant ; le serveur acquitte la r\'{e}ception de chaque trame. C'est la m\'{e}thode principale d'ingestion.

  \item \textbf{\texttt{WatchVehicle}} : streaming serveur (server-side streaming). Permet \`{a} un tableau de bord de s'abonner en temps r\'{e}el aux mouvements d'un v\'{e}hicule sp\'{e}cifique, via les Observables RxJS de NestJS.
\end{itemize}

\subsection{D\'{e}tection de g\'{e}ofencing}

Des zones g\'{e}ographiques (\textit{geofences}) sont d\'{e}finissables via l'API REST du service (\texttt{/api/geofences}). \`{A} chaque nouvelle position re\c{c}ue, le service calcule si le v\'{e}hicule se trouve en dehors d'une zone configur\'{e}e. En cas de franchissement, un \'{e}v\'{e}nement \texttt{GEOFENCING\_BREACH} est publi\'{e} sur le topic Kafka \texttt{flotte.location.events}, qui est imm\'{e}diatement trait\'{e} par le service \'{e}v\'{e}nements pour cr\'{e}er une alerte.

\subsection{Simulateur GPS}

Afin de valider l'architecture sans v\'{e}hicule physique, un \textbf{simulateur} a \'{e}t\'{e} d\'{e}velopp\'{e} en parall\`{e}le du service. Il s'agit d'un script Node.js autonome qui joue le r\^{o}le d'un client gRPC externe : il g\'{e}n\`{e}re un itin\'{e}raire factice (d\'{e}part de Rouen) et envoie des coordonn\'{e}es GPS r\'{e}alistes toutes les 3 secondes via \texttt{StreamPositions}.

Cette s\'{e}paration physique (le simulateur est un processus distinct, pas un module interne) garantit que la validation se fait par le r\'{e}seau, comme en production. Le simulateur est int\'{e}gr\'{e} \`{a} la stack Docker Compose et peut \'{e}galement servir de test de charge sur la base de donn\'{e}es temporelle.

\subsection{D\'{e}ploiement}

\begin{itemize}[leftmargin=*]
  \item \textbf{Docker Compose} : deux conteneurs ajout\'{e}s --- \texttt{location} (NestJS, port 3000 HTTP + 50051 gRPC) et \texttt{location-simulator} (client gRPC autonome).
  \item \textbf{Kubernetes} : Deployment et Service appliqu\'{e}s manuellement (manifestes d\'{e}di\'{e}s, hors chart \texttt{fleet-app} qui est con\c{c}u pour les services Java). La base \texttt{location\_db} est cr\'{e}\'{e}e dans le script \texttt{initdb} du chart \texttt{fleet-infra}. Les routes \texttt{/api/geofences} et \texttt{/api/locations} sont expos\'{e}es via l'Ingress Nginx.
\end{itemize}

% ============================================================
\newpage
\section{Tests Bruno}
% ============================================================

\begin{infobox}
\textbf{Contexte :} La collection Bruno couvre l'int\'{e}gralit\'{e} des services d\'{e}ploy\'{e}s. Cette semaine, 18 nouvelles requ\^{e}tes ont \'{e}t\'{e} ajout\'{e}es pour les deux nouveaux services : 9~REST et 6~GraphQL pour le service \'{e}v\'{e}nements, 3~REST pour le service localisation.
\end{infobox}

\begin{table}[H]
\centering
\small
\begin{tabular}{|l|c|c|}
  \hline
  \rowcolor{primary!15}
  \textbf{Service} & \textbf{Requ\^{e}tes REST} & \textbf{Requ\^{e}tes GraphQL} \\
  \hline
  service-vehicule    & 5 & 6 \\
  service-driver      & 8 & 4 \\
  service-maintenance & 5 & 6 \\
  service-events      & 9 \textit{(nouveau)} & 6 \textit{(nouveau)} \\
  service-location    & 3 \textit{(nouveau)} & --- \\
  \hline
  \textbf{Total}      & \textbf{30} & \textbf{22} \\
  \hline
\end{tabular}
\caption{R\'{e}partition des 52 requ\^{e}tes de la collection Bruno}
\end{table}

L'ex\'{e}cution compl\`{e}te sur l'environnement Minikube donne les r\'{e}sultats suivants :

\begin{table}[H]
\centering
\small
\begin{tabular}{|l|r|}
  \hline
  \rowcolor{primary!15}
  \textbf{M\'{e}trique} & \textbf{R\'{e}sultat} \\
  \hline
  Statut global  & \textcolor{success}{\textbf{PASS}} \\
  Requ\^{e}tes  & 53 / 53 \\
  Tests          & 81 / 81 \\
  \hline
\end{tabular}
\caption{R\'{e}sultat de \texttt{bru run -{}-env-file environments/Minikube.json -r}}
\end{table}

% ============================================================
\newpage
\section{Bilan et prochaines \'{e}tapes}
% ============================================================

\subsection{R\'{e}capitulatif de la semaine}

\begin{table}[H]
\centering
\small
\begin{tabular}{|l|l|c|}
  \hline
  \rowcolor{primary!15}
  \textbf{Fonctionnalit\'{e}} & \textbf{Service} & \textbf{Statut} \\
  \hline
  Consommation Kafka (maintenance + conducteurs)  & \'{E}v\'{e}nements & \textcolor{success}{\textbf{OK}} \\
  Consumer Kafka pr\^{e}t pour localisation       & \'{E}v\'{e}nements & \textcolor{success}{\textbf{OK}} \\
  API REST gestion des alertes                    & \'{E}v\'{e}nements & \textcolor{success}{\textbf{OK}} \\
  Int\'{e}gration GraphQL (queries + mutations)   & \'{E}v\'{e}nements & \textcolor{success}{\textbf{OK}} \\
  Tests unitaires (couverture $>$ 80\%)           & \'{E}v\'{e}nements & \textcolor{success}{\textbf{OK}} \\
  Ingestion GPS via gRPC                          & Localisation      & \textcolor{success}{\textbf{OK}} \\
  Stockage TimescaleDB (s\'{e}ries temporelles)   & Localisation      & \textcolor{success}{\textbf{OK}} \\
  D\'{e}tection g\'{e}ofencing + publication Kafka & Localisation     & \textcolor{success}{\textbf{OK}} \\
  API REST g\'{e}ofences                          & Localisation      & \textcolor{success}{\textbf{OK}} \\
  Simulateur GPS (Docker Compose)                 & Localisation      & \textcolor{success}{\textbf{OK}} \\
  D\'{e}ploiement Minikube (deux services)        & Infrastructure    & \textcolor{success}{\textbf{OK}} \\
  Tests Bruno : 53 requ\^{e}tes, 81/81 tests      & Infrastructure    & \textcolor{success}{\textbf{OK}} \\
  Subscriptions WebSocket (alertes temps r\'{e}el) & \'{E}v\'{e}nements & \textcolor{warning}{\textbf{S7}} \\
  \hline
\end{tabular}
\caption{Bilan de la semaine 5}
\end{table}

\subsection{Semaine 6}

La semaine~6 sera consacr\'{e}e au d\'{e}veloppement du \textbf{front-end} : mise en place de l'interface utilisateur permettant aux op\'{e}rateurs de visualiser et g\'{e}rer la flotte (v\'{e}hicules, conducteurs, alertes, carte de localisation).

\end{document}
